\documentclass[12pt]{article}
\usepackage{amsmath,amssymb,mathtools,bm,mathrsfs,extarrows,graphicx}
\usepackage[nointegrals]{wasysym}
\usepackage{pifont,stmaryrd}
\usepackage[margin=1in]{geometry}
\usepackage{hyperref}
\newcommand{\R}{\mathbb{R}}
\def\myparen#1{\left(#1\right)}
\DeclareMathOperator{\rank}{rank}
\providecommand{\reallywidehat}[1]{\mathpalette\PasTeXWideHat{#1}}
\providecommand{\PasTeXWideHat}[2]{%
  \begingroup\sbox0{$#1#2$}%
  \mathord{\overset{\resizebox{\wd0}{0.45ex}{$\wedge$}}{#1#2}}%
  \endgroup}
\begin{document}
\section*{Mathematics in PasTeX}
This reference covers symbols, structures, fonts, equation layouts, and custom
commands. Click a formula to edit it; the Equation tab provides insertion tools.
Press Space to leave a script, and use the arrow keys to move between fields.
See the user's guide for tables, figures, theorems, and bibliography editing.

\section{Inline, displayed, and numbered formulas}
Inline mathematics, such as $f\colon\R\to\R$ and $E=mc^2$, follows the surrounding
text. A displayed formula occupies its own line:
\[
 (a+b)^2=a^2+2ab+b^2.
\]
Select a formula and press Alt+= to turn it into body text. Use the layout
controls to change an inline formula into a display.

\begin{equation}\label{eq:gaussian}
 \int_{-\infty}^{\infty}e^{-x^2}\,dx=\sqrt{\pi}.
\end{equation}
Equation~\eqref{eq:gaussian} demonstrates a label and a clickable reference.
A custom tag replaces the automatic equation number:
\begin{equation}
 a^2+b^2=c^2.\tag{P}\label{eq:pythagoras}
\end{equation}
Compare Equation~\eqref{eq:pythagoras} with Equation~\eqref{eq:gaussian}.

\section{Symbols, relations, and arrows}
Common lowercase Greek letters and their variants:
\[
 \alpha,\,\beta,\,\gamma,\,\delta,\,\epsilon,\,\varepsilon,\,\zeta,\,\eta,
 \theta,\,\vartheta,\,\iota,\,\kappa,\,\lambda,\,\mu,\,\nu,\,\xi,\,\pi,
 \varpi,\,\rho,\,\varrho,\,\sigma,\,\varsigma,\,\tau,\,\upsilon,\,\phi,
 \varphi,\,\chi,\,\psi,\,\omega.
\]
Uppercase Greek letters:
\[
 \Gamma,\,\Delta,\,\Theta,\,\Lambda,\,\Xi,\,\Pi,\,\Sigma,\,\Upsilon,\,\Phi,
 \Psi,\,\Omega.
\]
Operators, relations, sets, and miscellaneous symbols:
\[
 a\pm b,\quad a\mp b,\quad a\times b,\quad a\div b,\quad a\cdot b,\quad
 A\cap B,\quad A\cup B,\quad A\setminus B,\quad A\oplus B,\quad A\otimes B
\]
\[
 x\leq y,\quad x\geq y,\quad x\neq y,\quad x\approx y,\quad x\equiv y,\quad
 x\sim y,\quad x\in A,\quad A\subseteq B,\quad \forall x,\quad \exists y,\quad
 \varnothing,\quad \infty,\quad \nabla,\quad \partial.
\]
\[
 A\to B,\quad A\leftarrow B,\quad A\leftrightarrow B,\quad
 A\Rightarrow B,\quad A\Leftarrow B,\quad A\Leftrightarrow B,\quad
 x\mapsto x^2,\quad A\hookrightarrow B,\quad A\xrightarrow{f}B.
\]

\subsection*{Installed symbol packages}
Equation, Symbols opens a searchable browser. Choose an installed package, or
the Comprehensive symbol list, then insert a command and its required package.
These examples use wasysym, pifont and stmaryrd:
\[
 \smiley\quad\text{\phone}\quad\CheckedBox\quad\HF\quad\LEFTcircle\quad
 \text{\ding{172}}\quad\llbracket x\rrbracket\quad a\apprle b.
\]
Double brackets also support fixed sizes and automatic growth:
\[
 \Bigl\llbracket x\Bigr\rrbracket\quad
 \left\llbracket\frac{x+1}{y}\right\rrbracket\quad
 \llfloor x\rrfloor\quad\llceil y\rrceil\quad
 \biginterleave_{i=1}^{n} A_i.
\]
Package-defined symbols use the installed font outlines and metrics. Some
catalogue entries require packages or programmable definitions beyond native
preview; use permanent preview to check their complete output.
Text symbols such as \verb|\phone| need \verb|\text{\phone}| inside mathematics.

\section{Fractions, roots, and scripts}
Compare inline $x_i^2$, $x^{n^2}_{i,j}$, $\sqrt{x}$, $\sqrt[n]{x+1}$,
and $\frac{1}{n^2}$. Empty arguments use editable placeholders in instant
editing; placeholders do not become part of the source.

\[
 \frac{a+b}{c+d}\qquad
 \frac{x}{\lVert\frac{x}{y}\rVert}\qquad
 \sqrt[3]{\frac{a+b}{c+d}}\qquad
 \binom{n}{k}
\]
Explicit text and display styles:
\[
 \tfrac{1}{2}+\dfrac{1}{2}\qquad
 \tbinom{n}{k}+\dbinom{n}{k}\qquad
 \sqrt{\pi}+\sqrt[n]{x+1}.
\]

\section{Operators and limits}
Named functions stay upright: $\sin x+\cos x$, $\log x$, $\exp x$,
$\max(a,b)$, and $\rank A$.
\[
 \sum_{n=0}^{\infty}\frac{x^n}{n!}=\mathrm{e}^x,\qquad
 \prod_{k=1}^{n}k=n!,\qquad
 \lim_{x\to0}\frac{\sin x}{x}=1
\]
\[
 \int_0^1x^2\,\mathrm{d}x=\frac{1}{3},\qquad
 \iint_D f(x,y)\,\mathrm{d}x\,\mathrm{d}y,\qquad
 \oint_C\bm{F}\cdot\mathrm{d}\bm{r}
\]
Inline limits, $\sum_{n=0}^{\infty}a_n$ and $\int_{\mathbb{R}}f(x)\,dx$,
use smaller operators than the displayed versions.
\[
 \bigcup_{i=1}^n A_i,\qquad
 \bigcap_{i=1}^n A_i,\qquad
 \bigoplus_{i=1}^n V_i,\qquad
 \bigotimes_{i=1}^n V_i.
\]

\subsection{Modular arithmetic}
The AMS commands distinguish a binary remainder operator from a congruence
qualifier. Their arguments remain editable in instant preview:
\[
 a\bmod n,\qquad a\equiv b\pmod{n},\qquad
 a\equiv b\mod{n},\qquad a\equiv b\pod{n}.
\]
Compare inline $a\equiv b\pmod{n}$ and the subscript $x_{a\bmod n}$ with
the display above. Named logarithms stay upright: $\lg(100)=2$.

\subsection{Math styles and invisible spacing}
The four arguments of \verb|\mathchoice| select display, text, script and
scriptscript forms. This command adapts as it enters smaller scripts:
\newcommand{\stylemark}{\mathchoice{D}{T}{S}{Q}}
Inline $\stylemark$ and $x_{\stylemark_{\stylemark}}$ compare with
\[
 \stylemark,\qquad a+\phantom{b}=c,\qquad a+\hphantom{b}=c.
\]
The phantom commands reserve the dimensions of their arguments.
The following fraction uses \verb|\mathpalette| to pass the current math
style to a two-argument helper.
\newcommand{\adaptfrac}[2]{{#1\frac{#2}{1+#2}}}
\newcommand{\stylefrac}[1]{\mathpalette\adaptfrac{#1}}
Compare inline $\stylefrac{x}$ and a script $a_{\stylefrac{x}}$ with
the displayed form
\[
  \stylefrac{x}.
\]

\subsection{Labelled arrows and functions}
Use Equation, Arrows for extensible arrows; edit the upper label in braces and
an optional lower label in brackets. Functions and limits have a separate gallery.
Insert math command opens a name-and-arguments editor for other commands.
\[
 A\xrightarrow[g]{f}B,\qquad
 A\xleftarrow{p}B,\qquad
 a\xlongequal[\text{definition}]{n\to\infty}b,\qquad
 U\xhookrightarrow{i}V\xRightarrow{T}W.
\]
\[
 \sin x+\cos x,\qquad \lim_{x\to0}\frac{\sin x}{x}=1,
 \qquad \sup_{x\in[0,1]}x=1,\qquad \limsup_{n\to\infty}a_n.
\]
\verb|\xlongequal| comes from \texttt{extarrows}; additional labelled arrows
come from \texttt{mathtools}. The preamble defines \verb|\reallywidehat| using
\texttt{amsmath} and \texttt{graphicx} to stretch over a long expression:
\[
 \widehat{abc}\qquad \reallywidehat{a+b+c+d+e+f}.
\]

\section{Matrices and determinants}
Choose the matrix type and dimensions from the Equation tab. Within a matrix,
arrow keys move through cells; the grid tools add or delete rows and columns.
\[
 \begin{matrix}a&b\\c&d\end{matrix}\quad
 \begin{pmatrix}1&2&3\\4&5&6\end{pmatrix}\quad
 \begin{bmatrix}a&b\\c&d\end{bmatrix}\quad
 \begin{Bmatrix}a&b\\c&d\end{Bmatrix}
\]
\[
 \begin{vmatrix}a&b\\c&d\end{vmatrix}=ad-bc,\qquad
 \begin{Vmatrix}x&0\\0&y\end{Vmatrix}.
\]
A small matrix can appear inline:
$A=\left(\begin{smallmatrix}1&0\\0&1\end{smallmatrix}\right)$.
An array allows explicit column alignment:
\[
 \begin{array}{rcl}
  x+y&=&3\\
  x-y&=&1
 \end{array}
\]

\section{Delimiters and spacing}
Automatically sized delimiters:
\[
 \left(\frac{1}{1+x^2}\right)\quad
 \left[\sum_{n=1}^N n\right]\quad
 \left\{\frac{a}{b},\sqrt[n]{c}\right\}\quad
 \left\langle\frac{x}{y},z\right\rangle
\]
\[
 \left\lvert\frac{x}{y}\right\rvert\quad
 \left\lVert\frac{x}{y}\right\rVert\quad
 \left.\frac{x^3}{3}\right|_0^1.
\]
Fixed sizes and unmatched sides:
\[
 \bigl(x\bigr)\quad \Bigl[x\Bigr]\quad \biggl\{x\biggr\}\quad
 \Biggl\langle x\Biggr\rangle\quad
 \left\{x\in\R\mid x>0\right..
\]
Click a space marker during instant editing to choose a preset or a width.
The following examples compare ordinary, thin, medium, thick, quad, and
double-quad spacing:
\[
 ab\qquad a\,b\qquad a\:b\qquad a\;b\qquad
 a\quad b\qquad a\qquad b.
\]

\section{Accents and annotations}
Compare fixed $\hat{x}$ and $\hat{xxxxxxxxxxxx}$ with extensible
$\widehat{x}$, $\widehat{xx}$, and $\widehat{xxxxxxxxxxxx}$.
\[
 \bar{x}+\overline{AB}+\vec{v}+\tilde{x}+\widetilde{xyz}
 +\dot{x}+\ddot{x}+\underline{a+b}
\]
\[
 \boxed{a+b}\qquad
 \overset{!}{=}\qquad \underset{x\to0}{\lim} f(x)\qquad
 \underbrace{a+\cdots+a}_{n\text{ terms}}.
\]

\section*{Mathematical alphabets}
Ordinary variables are italic; numbers and named operators are upright.
Text inside a formula uses the document's text font.
\[
 \begin{aligned}
  \text{Default}    &\quad ABCxyz0123456789+\alpha+\pi\\
  \text{Roman}      &\quad \mathrm{ABCxyz0123456789}\\
  \text{Italic}     &\quad \mathit{ABCxyz0123456789}\\
  \text{Bold roman} &\quad \mathbf{ABCxyz0123456789}\\
  \text{Sans serif} &\quad \mathsf{ABCxyz0123456789}\\
  \text{Typewriter} &\quad \mathtt{ABCxyz0123456789}
 \end{aligned}
\]
Special alphabets:
\[
 \begin{aligned}
  \text{Calligraphic}&\quad\mathcal{ABCDEFGHIJKLMNOPQRSTUVWXYZ}\\
  \text{Blackboard} &\quad\mathbb{ABCDEFGHIJKLMNOPQRSTUVWXYZ}\\
  \text{RSFS script}&\quad\mathscr{ABCDEFGHIJKLMNOPQRSTUVWXYZ}\\
  \text{Fraktur}    &\quad\mathfrak{ABCDEFxyz0123456789}
 \end{aligned}
\]
Use \verb|\boldsymbol| or \verb|\bm| for bold Greek letters and symbols:
\[
 \alpha+\beta+\pi+\sum_{n=0}^{\infty}a_n\qquad
 \boldsymbol{\alpha+\beta+\pi+\sum_{n=0}^{\infty}a_n}\qquad
 \mathbf{A}\bm{x}=\bm{b}.
\]
\[
 \mathscr{F}(f)(\xi)=
 \int_{\mathbb{R}}f(x)\,\mathrm{e}^{-2\pi\mathrm{i}x\xi}\,\mathrm{d}x,
 \qquad \mathcal{L}(V),\quad \mathfrak{g}.
\]

\section{Equation environments}
An alignment uses ampersands to mark matching columns:
\begin{align}
 (a+b)^2&=a^2+2ab+b^2\label{eq:expand}\\
 (a-b)^2&=a^2-2ab+b^2
\end{align}
Equation~\eqref{eq:expand} links to its numbered row.
\begin{align*}
 e^{i\theta}&=\cos\theta+i\sin\theta\\
 \nabla\cdot\mathbf{E}&=\frac{\rho}{\varepsilon_0}\\
 \sum_{n=1}^{\infty}\frac{1}{n^2}&=\frac{\pi^2}{6}
\end{align*}
Use aligned or gathered for smaller blocks inside a display:
\[
 \begin{aligned}
  x+y&=3\\
  x-y&=1
 \end{aligned}
 \qquad
 \begin{gathered}
  x=2\\
  x^2+y^2=5\\
  y=1
 \end{gathered}
\]
\begin{gather}
 a+b=b+a\\
 a(b+c)=ab+ac
\end{gather}
\begin{equation}
 \begin{split}
  (a+b+c)^2&=a^2+b^2+c^2\\
           &\quad+2ab+2bc+2ca
 \end{split}
\end{equation}
\[
 f(x)=\begin{cases}
  x^2&\text{if }x\geq0,\\
  -x&\text{otherwise.}
 \end{cases}
 \qquad a_{\text{initial}}=0.
\]

\section{Custom commands and literal source}
The preamble defines \verb|\R|, \verb|\myparen|, and \verb|\rank|:
\[
 \myparen{\frac{x+1}{2}}\in\R,\qquad \rank A=2.
\]
A definition in the body appears as an orange editable object. The definition
below contains bare mathematics, so its right-hand field uses math rendering:
\newcommand{\squareplus}[1]{#1^2+1}
\[
 \squareplus{x}+\squareplus{y}=x^2+y^2+2.
\]
A body-text definition uses the normal text font:
\newcommand{\greeting}[1]{Hello, \textit{#1}!}
\greeting{reader} The source remains a custom command.
Literal source is also editable: \verb|\sqrt[n]{x} + x^2|.
\end{document}
